\chapter{\methodologyname}
\label{chap:metodologia}

\section{Dise\~no general}

El mapeo se estructura en cuatro fases, cada una vinculada a un marco normativo de referencia (v\'ease tabla~\ref{tab:fases}).

\begin{table}[h]
\centering
\caption{Las cuatro fases de la metodolog\'ia}
\label{tab:fases}
\begin{tabularx}{\textwidth}{lXl}
\toprule
\textbf{Fase} & \textbf{Prop\'osito} & \textbf{Marco} \\
\midrule
Fase 1. Evidencia & Obtener requisitos del mundo real & IREB Elicitaci\'on \\
Fase 2. Formalizaci\'on & Requisitos con atributos normativos & IREB + ISO 29148 + INCOSE \\
Fase 3. Ejecuci\'on SpecKit & Observar el pipeline SDD & SpecKit SDD \\
Fase 4. An\'alisis del mapeo & Contrastar flujo observado vs. normativo & IREB + ISO 29148 + INCOSE \\
\bottomrule
\end{tabularx}
\end{table}

\section{Fase 1: Extracci\'on de evidencia}

La primera fase corresponde a la elicitaci\'on IREB. Se identifican cambios concretos en seis repositorios \emph{Open Source} representativos de distintos tipos de sistema (v\'ease tabla~\ref{tab:repos}).

\begin{table}[h]
\centering
\caption{Repositorios seleccionados}
\label{tab:repos}
\begin{tabularx}{\textwidth}{lX}
\toprule
\textbf{Repositorio} & \textbf{Tipo de sistema} \\
\midrule
Appwrite & Sistemas transaccionales (BaaS) \\
Medusa & Sistemas financieros (e-commerce) \\
Authentik & Identidad y seguridad \\
Cal.com & Sistemas colaborativos \\
Directus & Datos y contenido (headless CMS) \\
n8n & Integraci\'on y composici\'on de servicios \\
\bottomrule
\end{tabularx}
\end{table}

Para la selecci\'on de cambios se define una r\'ubrica con cuatro dimensiones (v\'ease tabla~\ref{tab:rubrica-seleccion}). Se descart\'o una r\'ubrica basada en IEEE 830 porque su aplicaci\'on sobre \emph{changelogs} generaba un rechazo pr\'acticamente universal de entradas, incluso aquellas con alto valor de trazabilidad.

\begin{table}[h]
\centering
\caption{R\'ubrica de selecci\'on de cambios (Fase 1)}
\label{tab:rubrica-seleccion}
\begin{tabularx}{\textwidth}{lX}
\toprule
\textbf{Dimensi\'on} & \textbf{Descripci\'on} \\
\midrule
Trazabilidad & La entrada referencia una PR o commit concreto con diff localizable \\
Se\'nales estructurales & Menciona superficie t\'ecnica, condici\'on o actor relevante \\
Utilidad para derivaci\'on & Describe comportamiento formalizable como requisito verificable \\
Defendibilidad & La inclusi\'on puede justificarse con razonamiento expl\'icito \\
\bottomrule
\end{tabularx}
\end{table}

Se almacenan \'unicamente textos literales con trazabilidad m\'inima (versi\'on, URL, PR, commit, se\'nales expl\'icitas), sin redactar requisitos en esta fase. La conservaci\'on del texto literal separa la observaci\'on de la formalizaci\'on y mejora la auditabilidad, en l\'inea con la distinci\'on IREB entre elicitaci\'on y documentaci\'on.

\section{Fase 2: Formalizaci\'on de requisitos}

La evidencia se transforma en requisitos siguiendo una plantilla con 12 atributos basados en IREB, ISO 29148 e INCOSE: ID, Nombre, Tipo, Descripci\'on formal, Fuente, Rationale, Prioridad, Verificaci\'on, Estado, Versi\'on, Dependencias y M\'odulo.

La descripci\'on formal debe cumplir ocho reglas obligatorias: expresar una \'unica obligaci\'on, comenzar por la f\'ormula ``El sistema deber\'a...'', contener un verbo observable, evitar t\'erminos subjetivos, incluir condiciones cuantificables, ser verificable, ser autosuficiente y no describir detalles de implementaci\'on.

Cada requisito debe superar un \emph{quality gate} de 12 criterios organizados en tres bloques (v\'ease Anexo~\ref{anexo:rubrica}):

\begin{description}
    \item[Bloque A. Identificaci\'on y gesti\'on (R1--R4c):] R1 (identificador \'unico), R2 (clasificaci\'on v\'alida), R3 (fuente trazable), R4 (prioridad asignada), R4b (dependencias documentadas), R4c (m\'odulo asignado).
    \item[Bloque B. Redacci\'on y verificabilidad (R5--R9):] R5 (estructura formal ``El sistema deber\'a...''), R6 (obligaci\'on \'unica), R7 (verbo observable), R8 (ausencia de ambig\"uedad cr\'itica), R9 (verificabilidad observable).
    \item[Bloque C. Completitud contextual (R10):] R10 (autosuficiencia del requisito).
\end{description}

Tres criterios son \textbf{bloqueantes}: R5 (estructura formal ausente), R7 (verbo no observable) y R9 (sin verificabilidad). Su incumplimiento provoca rechazo autom\'atico independientemente de la puntuaci\'on total.

\section{Fase 2.5: Derivaci\'on de MRS (Minimal Requirement Seeds)}

La Fase 2.5 introduce un paso intermedio entre la formalizaci\'on y la ejecuci\'on del pipeline. Su objetivo es transformar cada requisito formalizado IREB en un \emph{Minimal Requirement Seed} (MRS): un prompt m\'inimo y replicable que un usuario de \textsc{SpecKit} necesitar\'ia formular para que el agente derive de forma aut\'onoma el requisito formal equivalente.

Un MRS debe satisfacer tres propiedades: \textbf{suficiencia} (contiene toda la informaci\'on necesaria para inferir los campos obligatorios de la plantilla), \textbf{minimalidad} (no sobredetermina la soluci\'on) y \textbf{replicabilidad} (dos agentes independientes producen requisitos equivalentes al leer el mismo MRS). La estructura can\'onica sigue el formato:

\begin{quote}
\textit{Como [ROL], quiero que el sistema [COMPORTAMIENTO OBSERVABLE] [contexto], para [MOTIVACI\'ON].}
\end{quote}

Se derivaron 52 MRS a partir de los 52 requisitos formalizados, utilizando un script asistido por LLM con reglas deterministas documentadas. Cada MRS conserva la trazabilidad hacia su \texttt{REQ-*.md} de origen mediante el ID.

\section{Fase 3: Ejecuci\'on comparativa de SpecKit}

La Fase 3 constituye el n\'ucleo experimental del trabajo. Para cada caso del corpus se ejecutan dos flujos paralelos sobre el mismo repositorio en el \emph{commit} inmediatamente anterior al cambio de referencia, cambiando \'unicamente el tratamiento del requisito de entrada:

\subsection{Flow A (baseline) --- SpecKit sin marco IREB}

Pipeline m\'inimo de 4 pasos: \texttt{/speckit.specify} (recibiendo el texto literal del \emph{changelog}) $\rightarrow$ \texttt{/speckit.plan} $\rightarrow$ \texttt{/speckit.tasks} $\rightarrow$ \texttt{/speckit.implement}. Sin \texttt{constitution} personalizado, sin \texttt{clarify}, sin \texttt{checklist}, sin \texttt{analyze}.

\subsection{Flow B (IREB-enhanced) --- SpecKit con IREB Kit}

Pipeline ampliado de 9 pasos: \texttt{/speckit.constitution} (con 9 principios IREB personalizados) $\rightarrow$ \texttt{/speckit.specify} (recibiendo el REQ formalizado con 12 atributos) $\rightarrow$ \texttt{/speckit.clarify} $\rightarrow$ \texttt{/speckit.checklist} $\rightarrow$ \texttt{/speckit.plan} $\rightarrow$ \texttt{/speckit.tasks} $\rightarrow$ \texttt{/speckit.analyze} (pre-implementaci\'on) $\rightarrow$ \texttt{/speckit.implement} $\rightarrow$ \texttt{/speckit.analyze} (post-implementaci\'on).

Se desarrollaron dos versiones del IREB Kit:

\begin{description}
    \item[Kit v1:] Constitution + plantillas de spec, plan, checklist, clarify y analyze. Pipeline de 9 pasos sin mecanismo de control de alcance en la implementaci\'on.
    \item[Kit v2:] A\~nade \texttt{AGENTS.md} con 5 reglas universales (R0: no alucinar contexto, R1: trazabilidad, R2: verificabilidad, R3: minimalidad, R4: no ambig\"uedad) y un \textbf{scope contract} pre-implementaci\'on que acota expl\'icitamente los archivos a modificar. Este scope contract se genera antes del paso \texttt{implement} para prevenir \emph{scope creep}.
\end{description}

De los 52 casos del corpus, se seleccionaron 6 (2 por repositorio) de los 3 repositorios con preparaci\'on m\'as eficiente (n8n, Cal.com, Directus) para la ejecuci\'on comparativa en ambos flujos. Esta reducci\'on fue necesaria porque cada caso requiere dos ejecuciones manuales en Copilot Chat.

\section{Fase 4: Evaluaci\'on con r\'ubrica SRCI v3}

La Fase 4 aplica la r\'ubrica SRCI v3 (\emph{Structured Requirement Conformance Inspection}) sobre los artefactos generados por cada flujo. La r\'ubrica organiza 24 criterios en cuatro bloques que siguen el flujo del pipeline (v\'ease Anexo~\ref{anexo:rubrica-eval}):

\begin{description}
    \item[Bloque A (12 criterios, 35\%).] Preservaci\'on del requisito en \texttt{spec.md}. Eval\'ua si \textsc{SpecKit} conserva los 12 atributos IREB del requisito de entrada o los degrada. Aplica la misma r\'ubrica de validaci\'on de la Fase 2.
    \item[Bloque B (5 criterios, 30\%).] Conformidad de la implementaci\'on. Eval\'ua cobertura de la intenci\'on principal, proporcionalidad del cambio, ausencia de regresiones, correspondencia con el dominio t\'ecnico y evidencia de intenci\'on en el c\'odigo.
    \item[Bloque C (3 criterios, 20\%).] Verificaci\'on y testing. Eval\'ua presencia de tests, cobertura de los criterios de verificaci\'on y adecuaci\'on del tipo de test.
    \item[Bloque D (4 criterios, 15\%).] Trazabilidad del pipeline. Eval\'ua si los artefactos intermedios (\texttt{plan.md}, \texttt{tasks.md}, \texttt{checklist.md}) soportan la trazabilidad requisito $\rightarrow$ implementaci\'on.
\end{description}

Cuatro criterios son \textbf{no compensables} (su incumplimiento provoca clasificaci\'on autom\'atica como No conforme): A11 (alucinaci\'on de contexto), B1 (no cubre la intenci\'on principal), B3 (regresi\'on evidente) y C1 (sin tests).

La puntuaci\'on se calcula como suma ponderada de los cuatro bloques, normalizada a escala 0--10. Las categor\'ias finales son: Conforme ($\geq 8.0$), Parcialmente conforme (6.0--7.9), Conformidad baja (4.0--5.9) y No conforme ($<4.0$).